\chapter{Épreuves répétées et dénombrement}
\section*{Exercice 20.1 -- Trois épreuves}
Pour trois épreuves indépendantes de probabilité de succès $1/4$, calculer la probabilité d'exactement deux succès.

\section*{Exercice 20.2 -- Au moins une réussite}
Quatre essais indépendants réussissent chacun avec probabilité $0{,}2$. Calculer la probabilité d'au moins une réussite.

\section*{Exercice 20.3 -- Dénombrement}
Justifier qu'il y a $10$ façons de placer deux succès parmi cinq essais.

\section*{Exercice 20.4 -- Sans remise}
Une urne contient deux rouges et une bleue. On tire deux boules sans remise. Calculer $P(R_1\cap R_2)$ et comparer à $P(R_1)P(R_2)$.

\section*{Exercice 20.5 -- Indépendance deux à deux}
Sur deux bits équiprobables indépendants, considérer A : premier bit $1$, B : second bit $1$, C : bits différents. Vérifier les indépendances par paires et l'échec pour les trois.

\section*{Exercice 20.6 -- Marche aléatoire}
Après quatre pas $\pm1$ indépendants équiprobables, quelle est la probabilité de revenir à l'origine ? Et d'être en $2$ ?
